\subsection{All relative alternating forms in degrees nine and twenty-seven}

\begin{statement}{A.113}{Proved}
{The complete finite pairing spaces and simultaneous correction at three}
Let $a\in\Q^\times$ have $v_3(a)=1$, let $n\in\{2,3\}$,
and put
\[
 N=3^n,\qquad d=N/3,\qquad
 K_n=\Q(\alpha),\quad \alpha^N=a,\qquad
 K_{n-1}=\Q(\alpha^3),\qquad L_n=K_n(\mu_N).
 \tag{A.113.1}
\]
The relative norm-one torus and its actual three-torsion module are
\[
 \begin{gathered}
 T_n^{\rm rel}=\ker\bigl(\operatorname{Res}_{K_n/\Q}\mathbf G_m
       \xrightarrow{N}\operatorname{Res}_{K_{n-1}/\Q}\mathbf G_m\bigr),\\
 M_n^{\rm rel}=T_n^{\rm rel}[3]
       =\ker(P_N\longrightarrow P_d)\otimes\mu_3,
 \end{gathered}
 \tag{A.113.2}
\]
where $P_b=\F_3[\Z/b\Z]$ and the displayed map sums coordinates
congruent modulo $d$.
Let $G=\operatorname{Gal}(L_n/\Q)=C_N\rtimes(\Z/N\Z)^\times$.
Then
\[
 \dim_{\F_3}\operatorname{Hom}_G
       (\mathop{\bigwedge}\nolimits^2M_n^{\rm rel},\mu_3)=n.
 \tag{A.113.3}
\]

Here is an explicit basis of this entire space.
Write $\F_3[C_N]=\F_3[Z]/(Z^N)$, with $Z=x-1$, so that
$M_n^{\rm rel}=Z^d\F_3[Z]/(Z^N)\otimes\mu_3$.
On each block of three coordinates use its sum-zero relation and set
\[
 \mathcal B_*(v,w)=\sum_{h=0}^{d-1}
                 (v_hw_{h+d}-v_{h+d}w_h),
       \qquad z=\frac{Z^2}{1+Z}.
 \tag{A.113.4}
\]
The expression is a $G$-equivariant nondegenerate alternating form
with values in $\mu_3$.
For $0\le j<n$ put
\[
 \mathcal B_j(v,w)=\mathcal B_*(v,z^{\,d-3^j}w).
 \tag{A.113.5}
\]
These are the promised basis, with
\[
 \operatorname{rank}\mathcal B_j=2\cdot3^j,\qquad
 \operatorname{rad}\mathcal B_j
       =Z^{d+2\cdot3^j}\F_3[Z]/(Z^N)\otimes\mu_3.
 \tag{A.113.6}
\]
The projective ranks are respectively
$\{6:3,\ 2:1\}$ for $n=2$, and
$\{18:9,\ 6:3,\ 2:1\}$ for $n=3$, where each value after a
colon is the number of projective forms of that rank.

Identify $\tfrac13\Z/\Z$ with $\F_3$.
The simultaneous, explicitly defined local map
\[
 \begin{split}
 T_n^{\rm rel}(\Q_3)/T_n^{\rm rel}(\Q_3)^3
      &\longrightarrow\F_3^n,\\
 u&\longmapsto
 \left(\operatorname{inv}_3\left(
       \tfrac12(\mathcal B_j)_*
          (\kappa_T(u)\smile\kappa_T(u))\right)\right)_{j=0}^{n-1}
 \end{split}
 \tag{A.113.7}
\]
is onto.  Consequently every collection of the quadratic
reciprocity values obtained from equivariant alternating forms
on this precise relative module can be compensated at three,
when that local multiplier is free.
This assertion supplies the local map (A.113.7); it does not
construct a global multiplier with prescribed new ramification.
\end{statement}

\begin{proof}
\emph{The norm map and the first form.}
Restriction of the embedding indexed by $i\bmod N$ to $K_{n-1}$
depends on $i\bmod d$.  The norm is the product on that block.
This proves the finite-torsion map in (A.113.2), including its
cyclotomic twist.  On the translation group algebra its kernel
is generated by $x^d-1=Z^d$, giving the stated polynomial module.
Each block sum-zero space has dimension two, and its determinant
form is the summand in (A.113.4).  Hence the full form is
nondegenerate.
A translation cyclically permutes a block when it crosses its
end; a three-cycle has determinant one on the sum-zero plane.
A unit $u\bmod N$ induces, within a block, an affine permutation
of three indices whose determinant is $u\bmod3$.
The two twists multiply to one.  Thus the sum in (A.113.4)
is equivariant with multiplier equal to the action on $\mu_3$.
This proves the exact coefficient map used in the statement.

\emph{Classification of all equivariant alternating forms.}
Put $D=2d=\dim M_n^{\rm rel}$.
Translation acts as the single Jordan block $1+Z$ on the basis
$Z^d,Z^{d+1},\ldots,Z^{N-1}$.
Its commuting endomorphisms are therefore exactly
the multiplication operators $f(Z)$ in $\F_3[Z]/(Z^D)$:
the image of the first cyclic vector determines the operator.
For any form $\mathcal B$ there is a unique operator $F$ with
\[
                \mathcal B(v,w)=\mathcal B_*(v,Fw).
 \tag{A.113.8}
\]
Equivariance of both forms is equivalent to $F$ commuting with
$G$.  Alternation of $\mathcal B$ is equivalent to $F$ being
self-adjoint for $\mathcal B_*$.

The unit group is generated by two at these two values of $N$.
Conjugation by its generator sends the operator $Z$ to
$2Z+Z^2$.
The adjoint of $Z$ is
$(1+Z)^{-1}-1=-Z/(1+Z)$, because translation is an isometry.
The latter substitution is also the action of the involution
$-1\in G$ on the operator algebra.
Thus every endomorphism commuting with $G$ is automatically
self-adjoint.
It remains to compute
\[
 \left\{f\in\F_3[Z]/(Z^{2d}):
                         f(2Z+Z^2)=f(Z)\right\}.
 \tag{A.113.9}
\]

The involution-fixed subring is exactly
$\F_3[z]/(z^d)$ for $z=Z^2/(1+Z)$.
Indeed $z$ is fixed, its powers have distinct leading terms
$Z^{2i}$, and the involution is diagonalizable with exactly
$d$ eigenvalues equal to one: its triangular diagonal on
$1,Z,\ldots,Z^{2d-1}$ is $1,-1,1,-1,\ldots$.
Conjugation by two induces the literal substitution
\[
                          z\longmapsto z+z^2.
 \tag{A.113.10}
\]
For $d=3$, coefficient comparison in
$f(z+z^2)-f(z)\bmod z^3$ gives
\[
                 f\in\langle1,z^2\rangle.
 \tag{A.113.11}
\]
For $d=9$, write $f=\sum_{i=0}^8c_i z^i$.
The coefficients of $z^2,z^3,z^5,z^6,z^7,z^8$ give, in order,
\[
 c_1=c_2=c_4=0,\qquad
 c_3+2c_5=0,\qquad c_5=0,\qquad c_7=0.
 \tag{A.113.12}
\]
There are no further conditions: $1,z^6,z^8$ are fixed
modulo $z^9$ by direct substitution.
Thus (A.113.9) is $\langle1,z^6,z^8\rangle$ in this case.
This proves completeness and linear independence of (A.113.5).

Multiplication by $z^{d-3^j}$ has leading power
$Z^{2d-2\cdot3^j}$ times a unit.
Its kernel on $M_n^{\rm rel}$ is exactly the ideal in
(A.113.6).  Nondegeneracy of $\mathcal B_*$ proves the
rank and radical formulas.  In a linear combination the
least nonzero power of $z$ determines the rank.
Counting its remaining free coefficients after projective
normalization gives the two displayed projective rank lists.

\emph{Local cohomology of the exact subquotients.}
The local group at three maps onto $G$ by the full local
Kummer degree and holomorph calculation of A.105.
For every submodule $Z^{d+s}\F_3[Z]/(Z^N)\otimes\mu_3$
with $s$ even, its translation coinvariant has involution
character $(-1)^{d+s}(-1)=1$.
It has no equivariant quotient $\mu_3$.
By local Tate duality,
\[
 H^2\bigl(\Q_3,Z^{d+s}\F_3[Z]/(Z^N)\otimes\mu_3\bigr)=0
                         \quad(s\text{ even}).
 \tag{A.113.13}
\]
The zero submodule satisfies the same assertion.
For a quotient between two such even-index submodules,
its translation-fixed line is its last power, of exponent
$d+t-1$ with $t$ even.  The involution acts there by $-1$.
Hence its $H^0$ is zero.
When the quotient has a nondegenerate alternating form into
$\mu_3$, it is its own Cartier dual, so its $H^2$ is also zero.
Local duality and the Euler characteristic
\cite[I, Corollary 2.3 and Theorem 2.8]{BIR_MilneADT}
then make its $H^1$ a vector space of dimension equal to the
coefficient dimension, with a nondegenerate symmetric
cup-product pairing.

\emph{Simultaneous values and the cross terms.}
First use $\mathcal B_0$.
Its quotient by its radical has dimension two, and
(A.113.13) makes the map on $H^1$ onto that quotient's
$H^1$.  Its half-diagonal quadratic form represents all of
$\F_3$, as in A.110.  Therefore the first value may be chosen.

For the next value, and then for the third when $n=3$, put
$h=3^{j-1}$ and restrict corrections to
\[
 R_{j-1}=\operatorname{rad}\mathcal B_{j-1}
            =Z^{d+2h}\F_3[Z]/(Z^N)\otimes\mu_3.
 \tag{A.113.14}
\]
Every earlier form vanishes when either argument belongs
to $R_{j-1}$.  Thus adding a cocycle with values in this
module changes none of the already chosen quadratic values,
including their cross terms.

The radical of the restriction of $\mathcal B_j$ to
$R_{j-1}$ is exactly
\[
 R'_{j-1}=Z^{d+4h}\F_3[Z]/(Z^N)\otimes\mu_3,
          \qquad \dim(R_{j-1}/R'_{j-1})=2h.
 \tag{A.113.15}
\]
To see this without a rank inference, in the cyclic module
of length $D$ the $\mathcal B_*$-orthogonal of $Z^sM$
is $Z^{D-s}M$: the adjoint of $Z$ is a unit times $Z$,
and the dimensions give equality.
The operator defining $\mathcal B_j$ has valuation
$D-6h$.  The condition that its value lie in the
orthogonal of $Z^{2h}M$ is therefore exactly membership
in $Z^{4h}M$, which proves (A.113.15).

By (A.113.13), $H^1(\Q_3,R_{j-1})$ maps onto the
$H^1$ of the quotient in (A.113.15).
Choose an $\F_3$-linear section of this finite-dimensional
surjection.  On its image the quadratic part coming from
$\mathcal B_j$ is nondegenerate, of dimension $2h\ge2$.
For an already chosen class $\gamma$, its value after adding
such a correction $\delta$ is
\[
 \mathcal Q_j(\gamma+\delta)
   =\mathcal Q_j(\gamma)+\mathcal Q_j(\delta)
       +\operatorname{inv}_3
                      (\mathcal B_j)_*(\gamma\smile\delta).
 \tag{A.113.16}
\]
The last term is a linear functional of $\delta$.
Every linear functional on a nondegenerate quadratic space
in odd characteristic is obtained by polar pairing with a
vector.  Completing the square therefore reduces
(A.113.16) to that same nondegenerate quadratic form plus
a constant.  A nondegenerate quadratic form of dimension
at least two over $\F_3$ represents every scalar.
The next value can consequently be chosen without changing
the earlier ones.  This proves simultaneous surjectivity
on $H^1(\Q_3,M_n^{\rm rel})$.

Finally the two local fields in (A.113.1) have unique
places above three, and their relative degree-three
norm is onto by A.105 (take $Q=1$).
Shapiro and Hilbert 90 identify
\[
 H^1(\Q_3,T_n^{\rm rel})
    =K_{n-1,3}^{\times}/N_{K_{n,3}/K_{n-1,3}}K_{n,3}^{\times}=0.
 \tag{A.113.17}
\]
The local Kummer map therefore identifies the source of
(A.113.7) with all of $H^1(\Q_3,M_n^{\rm rel})$.
This proves (A.113.7) with actual relative norm-one elements.
\end{proof}

The retained \path{code/relative_norm_forms_certificate.py},
its output, and its exact matrix data independently solve all
the equivariance equations in the two stated dimensions.
They record complete bases, all projective ranks and radicals.
The coefficient computation (A.113.9)--(A.113.12) proves
completeness without an extrapolation to higher $n$.
